%! Polynomial Long and Synthetic Division — Unit 3, Lesson 3.5 — Guided Notes
% THE in-class document, in the MAIN IDEAS / NOTES shape (LESSON_SHAPE.md §1): the vocabulary
% box, the hook, then ONE two-column guidednotes table. Each row is one idea — a label on the
% left; on the right one or two COMPLETE printed sentences, ONE large pre-drawn display, then a
% two-across grid of numbered problems. Problems run 1..18.
%   I DO   : the teacher reads the definition, marks up the display, works the first problem
%            of each grid (1, 3, 7, 9).
%   WE DO  : the class works the rest of each grid, students holding the pen, every one
%            cold-called.
%   YOU DO : the last row — problems 15–18, alone, for 17 minutes while the teacher circulates.
% The crux is row 4, FACTOR THEOREM: problems 11 and 12 divide x^3 - 7x + 6 by x - 2 WITHOUT and
% WITH the 0x^2 placeholder — the remainders disagree (-4 vs 0) — and problem 13 settles it with
% p(2) = 0. You Do 16 and 17 carry the missing power again; 15 and 18 carry the sign of c.
%
% DIVISION FRAMES: each long or synthetic division is a pre-drawn tabular whose columns have
% fixed widths and whose every row carries the same strut (\dstrut), so an empty cell in the
% blank and a filled \ans cell here are the same height. The blank prints the empty frame
% (divisor, dividend, rules); the key fills it. Nothing else in a frame differs.
%
% PAGE LOCKSTEP with notes/: the vocabulary write lines -> \ansline, the hook's lines -> one
% \ansline, the frame cells -> \ans, and the answer argument of every \pcell, \writespace and
% \labelbox filled. Nothing else differs.
% Inside a cell `\\` ENDS THE ROW — break lines with \par (`\\` inside tabular is safe).
% A row cannot break across pages: the figure and EACH grid row are separate sub-rows.
\documentclass[10pt]{article}
\usepackage{precalculus-article}
\usepackage{precalculus-boxes}
\newcommand{\TallMath}[1]{$\displaystyle #1\rule[-1.4em]{0pt}{3.2em}$}
% Division-frame columns: T = one term of a long division, K = one synthetic-division number.
\newcolumntype{T}{>{\centering\arraybackslash}p{0.98cm}}
\newcolumntype{K}{>{\centering\arraybackslash}p{0.62cm}}
\newcommand{\dstrut}{\rule[-3.5pt]{0pt}{12pt}}

\begin{document}

\pageheader{Unit 3, Lesson 3.5}{Guided Notes}

\vspace{0.12in}

\boxguard
\begin{vocabbox}
Fill in each term as we define it together.
\par\vspace{2pt}
\noindent\textbf{\textcolor{plum}{Division algorithm:}}\\[1pt]
\writeline\\[3pt]
\noindent\textbf{\textcolor{plum}{Placeholder:}}\\[1pt]
\writeline\\[3pt]
\noindent\textbf{\textcolor{plum}{Synthetic division:}}\\[1pt]
\writeline\\[3pt]
\noindent\textbf{\textcolor{plum}{Remainder Theorem:}}\\[1pt]
\writeline\\[3pt]
\noindent\textbf{\textcolor{plum}{Factor Theorem:}}\\[1pt]
\writeline
\end{vocabbox}

\vspace{0.1in}

\boxguard
\begin{hookbox}
None of the 3.4 patterns fits $p(x) = x^3 - 7x + 6$: no common factor, no grouping, not a sum
or difference of cubes. Maya says, ``Then there is no way to tell whether $x - 2$ is a factor.''
Jordan says, ``Just plug in 2.''

How could plugging in one number tell you anything about a factor?
\writelines{2}
\end{hookbox}

\small
\begin{guidednotes}
% ── ROW 1: long division — set up like numbers ───────────────────────────────
\mainidea[Like dividing numbers]{Long Division} &
Dividing a polynomial $f(x)$ by $g(x)$ gives a \textbf{quotient} $q(x)$ and a
\textbf{remainder} $r(x)$, with $f(x) = g(x)\,q(x) + r(x)$ and the degree of $r$ less than
the degree of $g$. Write both polynomials highest power first, with a $0$ for any missing power.

\vspace{3pt}
\stepnum{1}\ Divide the leading terms. \quad \stepnum{2}\ Multiply back. \quad
\stepnum{3}\ Subtract --- change every sign. \quad \stepnum{4}\ Bring down. Repeat until
the remainder's degree is less than the divisor's.
\\
& \begin{minipage}[c]{0.36\linewidth}\centering
{\setlength{\tabcolsep}{2.5pt}\begin{tabular}{@{}>{\dstrut}r@{\,}ccc@{}}
            &        & \textbf{\color{plum}1} & \textbf{\color{plum}3} \\ \cline{2-4}
 $12\,\big)$ & $1$    & $5$    & $7$    \\
            & $-(1$  & $2)$   &        \\ \cline{2-3}
            &        & $3$    & $7$    \\
            &        & $-(3$  & $6)$   \\ \cline{3-4}
            &        &        & \textbf{\color{redacc}1} \\
\end{tabular}}

\vspace{3pt}
$157 = 12 \cdot 13 + 1$
\end{minipage}\hfill
\begin{minipage}[c]{0.6\linewidth}\centering
{\setlength{\tabcolsep}{2.5pt}\begin{tabular}{@{}>{\dstrut}r@{\,}ccc@{}}
              &          & \textbf{\color{plum}$x$} & \textbf{\color{plum}$+3$} \\ \cline{2-4}
 $x+2\,\big)$ & $x^2$    & $+5x$    & $+7$    \\
              & $-(x^2$  & $+2x)$   &         \\ \cline{2-3}
              &          & $3x$     & $+7$    \\
              &          & $-(3x$   & $+6)$   \\ \cline{3-4}
              &          &          & \textbf{\color{redacc}$1$} \\
\end{tabular}}

\vspace{3pt}
$x^2 + 5x + 7 = (x + 2)\,\cdot$ \labelbox{1.6cm}{} $+$ \labelbox{0.9cm}{}
\end{minipage}
\\
& \notesprompt{Divide in the frame. Then write dividend $=$ divisor $\cdot$ quotient $+$ remainder.}
\begin{probgrid}{|Y|Y|}
\pcell{1}{$(2x^3 - 3x^2 - 5x + 6) \div (x - 2)$\par\smallskip
{\setlength{\tabcolsep}{1.5pt}\begin{tabular}{@{}>{\dstrut}r@{\,}TTTT@{}}
             &   &  &  &  \\ \cline{2-5}
 $x-2\,\big)$ & $2x^3$ & $-3x^2$ & $-5x$ & $+6$ \\
             &  &  & & \\ \cline{2-3}
             & &  &  & \\
             & &  &  & \\ \cline{3-4}
             & & &  &  \\
             & & &  &  \\ \cline{4-5}
             & & & &  \\
\end{tabular}}}{1.1cm}{} &
\pcell{2}{$(x^3 + 2x^2 - x + 3) \div (x^2 + 1)$. Why does it stop with an $x$ term left?\par\smallskip
{\setlength{\tabcolsep}{1.5pt}\begin{tabular}{@{}>{\dstrut}r@{\,}TTTT@{}}
               &   &   &  &  \\ \cline{2-5}
 $x^2+1\,\big)$ & $x^3$ & $+2x^2$ & $-x$ & $+3$ \\
               &  & &  & \\ \cline{2-4}
               & &  &  &  \\
               & &  & &  \\ \cline{3-5}
               & & &  &  \\
\end{tabular}}}{1.2cm}{} \\ \hline
\end{probgrid}
\\ \hline
% ── ROW 2: synthetic division — the same numbers, faster ─────────────────────
\mainidea[A shortcut for $x - c$]{Synthetic} &
When the divisor is $x - c$, \textbf{synthetic division} keeps only the coefficients. Write
$c$ in the corner and the coefficients in a row (a $0$ for any missing power). Bring down the
first; then multiply by $c$, add down the next column, and repeat. The last number is the
remainder; the others are the quotient's coefficients, one degree lower than the dividend.
\\
& \begin{minipage}[c]{0.56\linewidth}\centering
{\setlength{\tabcolsep}{1.5pt}\begin{tabular}{@{}>{\dstrut}r@{\,}TTTT@{}}
             &          & \textbf{\color{plum}$x^2$} & \textbf{\color{plum}$+x$} & \textbf{\color{plum}$-2$} \\ \cline{2-5}
 $x-3\,\big)$ & $x^3$    & $-2x^2$  & $-5x$    & $+6$   \\
             & $-(x^3$  & $-3x^2)$ &          &        \\ \cline{2-3}
             &          & $x^2$    & $-5x$    &        \\
             &          & $-(x^2$  & $-3x)$   &        \\ \cline{3-4}
             &          &          & $-2x$    & $+6$   \\
             &          &          & $-(-2x$  & $+6)$  \\ \cline{4-5}
             &          &          &          & \textbf{\color{redacc}$0$} \\
\end{tabular}}

\vspace{2pt}
{\footnotesize\itshape Long division of $(x^3 - 2x^2 - 5x + 6) \div (x - 3)$}
\end{minipage}\hfill
\begin{minipage}[c]{0.42\linewidth}\centering
{\setlength{\tabcolsep}{2pt}\begin{tabular}{@{}>{\dstrut}K|KKKK@{}}
 $3$ & $1$ & $-2$ & $-5$ & $6$ \\
     &     & $3$  & $3$  & $-6$ \\ \cline{2-5}
     & \textbf{\color{plum}$1$} & \textbf{\color{plum}$1$} & \textbf{\color{plum}$-2$} & \textbf{\color{redacc}$0$} \\
\end{tabular}}

\vspace{2pt}
{\footnotesize\itshape The same division, synthetic}

\vspace{6pt}
\raggedright
$1,\ 1,\ -2$ are the \labelbox{2.3cm}{} of the quotient $x^2 + x - 2$.\par
\vspace{3pt}
The last number, $0$, is the \labelbox{2.3cm}{}.
\end{minipage}
\\
& \notesprompt{Divide by synthetic division. Say what $c$ is first.}
\begin{probgrid}{|Y|Y|}
\pcell{3}{Redo problem 1 the short way: $(2x^3 - 3x^2 - 5x + 6) \div (x - 2)$.\par\smallskip
{\setlength{\tabcolsep}{2pt}\begin{tabular}{@{}>{\dstrut}K|KKKK@{}}
  &  &  &  &  \\
     &     &   &   &  \\ \cline{2-5}
     &  &  &  &  \\
\end{tabular}}}{1.1cm}{} &
\pcell{4}{$(x^3 + 5x^2 + 2x - 4) \div (x + 2)$\par\smallskip
{\setlength{\tabcolsep}{2pt}\begin{tabular}{@{}>{\dstrut}K|KKKK@{}}
  &  &  &  &  \\
     &     &   &   &  \\ \cline{2-5}
     &  &  &  &  \\
\end{tabular}}}{1.1cm}{} \\ \hline
\end{probgrid}
\\
& \begin{probgrid}*{|Y|Y|}
\pcell{5}{$(3x^3 - 4x^2 + 2x - 5) \div (x - 1)$\par\smallskip
{\setlength{\tabcolsep}{2pt}\begin{tabular}{@{}>{\dstrut}K|KKKK@{}}
  &  &  &  &  \\
     &     &   &   &  \\ \cline{2-5}
     &  &  &  &  \\
\end{tabular}}}{1.1cm}{} &
\pcell{6}{Synthetic division needs a divisor of the form $x - c$. Give $c$, or write ``long division.''\par\smallskip
(a) $x - 4$ \qquad (b) $x + 5$ \qquad (c) $x^2 + 1$ \qquad (d) $x$}{1.4cm}{} \\ \hline
\end{probgrid}
\\ \hline
% ── ROW 3: the Remainder Theorem — two roads to p(c) ─────────────────────────
\mainidea[Two roads to $p(c)$]{Remainder Theorem} &
\textbf{The Remainder Theorem:} when $p(x)$ is divided by $x - c$, the remainder is $p(c)$.
The reason is the division algorithm itself: $p(x) = (x - c)\,q(x) + r$, so at $x = c$ the
first term is worth nothing and only the remainder is left.
\\
& \begin{center}
\fcolorbox{plum}{white}{\begin{minipage}[c][2.3cm][c]{0.4\linewidth}\centering
\textbf{\color{plum}Road 1: substitute}\par\vspace{4pt}
$p(x) = x^3 - 2x^2 + 4$\par\vspace{3pt}
$p(3) = 3^3 - 2(3)^2 + 4$\par
$= 27 - 18 + 4 = \textbf{\color{redacc}13}$
\end{minipage}}
\quad\raisebox{0.9cm}{\Large$\Longleftrightarrow$}\quad
\fcolorbox{plum}{white}{\begin{minipage}[c][2.3cm][c]{0.4\linewidth}\centering
\textbf{\color{plum}Road 2: divide by $x - 3$}\par\vspace{2pt}
{\setlength{\tabcolsep}{2pt}\begin{tabular}{@{}>{\dstrut}K|KKKK@{}}
 $3$ & $1$ & $-2$ & $0$ & $4$ \\
     &     & $3$  & $3$ & $9$ \\ \cline{2-5}
     & $1$ & $1$  & $3$ & \textbf{\color{redacc}$13$} \\
\end{tabular}}
\end{minipage}}
\end{center}

\vspace{2pt}
At $x = 3$ the factor $(x - 3)$ equals \labelbox{0.9cm}{}, so $p(3)$ is exactly the
\labelbox{2.3cm}{}.
\\
& \notesprompt{Find each value. Two roads must land on the same number.}
\begin{probgrid}{|Y|Y|}
\pcell{7}{$p(x) = x^3 + 3x^2 - 4x + 1$. Find $p(-2)$ by substituting, then by synthetic division.\par\smallskip
{\setlength{\tabcolsep}{2pt}\begin{tabular}{@{}>{\dstrut}K|KKKK@{}}
  &  &  &  &  \\
     &     &   &   &  \\ \cline{2-5}
     &  &  &  &  \\
\end{tabular}}}{1.3cm}{} &
\pcell{8}{Without dividing, find the remainder when $2x^3 - 5x^2 + x + 4$ is divided by $x - 3$.\par\smallskip
Then: if a remainder is $0$, what does that say about $p(c)$?}{2.4cm}{} \\ \hline
\end{probgrid}
\\ \hline
% ── ROW 4: THE CRUX — the Factor Theorem, and the missing power ──────────────
\mainidea[Remainder zero]{Factor Theorem} &
\textbf{The Factor Theorem:} $(x - c)$ is a factor of $p(x)$ exactly when $p(c) = 0$. So a
remainder of $0$, a factor, a zero, and an $x$-intercept are four ways of saying one thing.
Once one zero is known, divide its factor out and factor the quotient (3.4) to find the rest.
\\
& \begin{minipage}[c]{0.5\linewidth}\centering
\begin{tikzpicture}[x=0.62cm, y=0.11cm, >=stealth]
  \draw[linegray, very thin, xstep=1, ystep=5] (-4,-15) grid (3,15);
  \draw[->, thick] (-4,0) -- (3.5,0) node[right] {\small $x$};
  \draw[->, thick] (0,-15) -- (0,16.5) node[above] {\small $p(x)$};
  \foreach \x in {-3,-2,-1,1,2,3} \node[below, font=\scriptsize] at (\x,0) {\x};
  \foreach \y in {-10,10} \node[left, font=\scriptsize] at (0,\y) {\y};
  \draw[very thick, plum, domain=-3.55:2.95, samples=80, smooth]
       plot (\x, {\x*\x*\x - 7*\x + 6});
  \foreach \x in {-3,1,2} \fill[plum] (\x,0) circle (2.3pt);
\end{tikzpicture}

{\footnotesize\itshape $p(x) = x^3 - 7x + 6$, the hook's polynomial}
\end{minipage}\hfill
\begin{minipage}[c]{0.47\linewidth}\centering
\fcolorbox{plum}{lilac}{\begin{minipage}{0.92\linewidth}\centering\small
remainder $0$ when $p(x) \div (x - c)$\par
$\Updownarrow$\par
$(x - c)$ is a factor: $p(x) = (x - c)\,q(x)$\par
$\Updownarrow$\par
$p(c) = 0$: \ $c$ is a zero\par
$\Updownarrow$\par
$(c, 0)$ is an $x$-intercept
\end{minipage}}

\vspace{6pt}
\raggedright
From the graph, $p(2) = $ \labelbox{0.9cm}{}, so $x - 2$ \labelbox{1.6cm}{} a factor.
\end{minipage}
\\
& \notesprompt{Divide out the factor you know; factor what is left.}
\begin{probgrid}{|Y|Y|}
\pcell{9}{Problems 1 and 3 divided $2x^3 - 3x^2 - 5x + 6$ by $x - 2$: quotient $2x^2 + x - 3$, remainder $0$. Write the cubic fully factored and list all three zeros.}{1.6cm}{} &
\pcell{10}{$-3$ is a zero of $x^3 + 2x^2 - 5x - 6$. Which factor does that give? Divide it out; find all the zeros.\par\smallskip
{\setlength{\tabcolsep}{2pt}\begin{tabular}{@{}>{\dstrut}K|KKKK@{}}
  &  &  &  &  \\
     &     &   &   &  \\ \cline{2-5}
     &  &  &  &  \\
\end{tabular}}}{1.3cm}{} \\ \hline
\end{probgrid}
\\
& \notesprompt{Is $x - 2$ a factor of $p(x) = x^3 - 7x + 6$? Ana and Ben each start a synthetic division.}
\begin{probgrid}*{|Y|Y|}
\pcell{11}{\textbf{Ana} writes the row $1 \quad -7 \quad 6$. Finish her division. What does it say?\par\smallskip
{\setlength{\tabcolsep}{2pt}\begin{tabular}{@{}>{\dstrut}K|KKK@{}}
 $2$ & $1$ & $-7$ & $6$ \\
     &     &   &  \\ \cline{2-4}
     &  &  &  \\
\end{tabular}}}{1.3cm}{} &
\pcell{12}{\textbf{Ben} writes the row $1 \quad 0 \quad -7 \quad 6$. Finish his division. What does it say?\par\smallskip
{\setlength{\tabcolsep}{2pt}\begin{tabular}{@{}>{\dstrut}K|KKKK@{}}
 $2$ & $1$ & $0$ & $-7$ & $6$ \\
     &     &   &   &  \\ \cline{2-5}
     &  &  &  &  \\
\end{tabular}}}{1.3cm}{} \\ \hline
\end{probgrid}
\\
& \begin{probgrid}*{|Y|Y|}
\pcell{13}{Settle it: find $p(2)$ by substituting. Whose remainder does the Remainder Theorem back? Then factor Ben's quotient and list all three zeros --- check them on the graph.}{2.2cm}{} &
\pcell{14}{Multiply out Ana's answer, $(x - 2)(x - 5) - 4$. Which polynomial did she \emph{really} divide? Which column did her $-7$ land in?}{2.2cm}{} \\ \hline
\end{probgrid}
\\
& \fcolorbox{plum}{lilac}{\parbox{\dimexpr\linewidth-2\fboxsep-2\fboxrule\relax}{\centering
\textbf{A missing power still gets a column. Write $0$ for it --- in long division and in
synthetic division --- or every number after it slides one column too far left, and the
division answers a different question.}}}

\vspace{3pt}
\textbf{In your own words:} why did Ana's remainder disagree with $p(2)$?
\writespace{1.2cm}{}
\\ \hline
% ── LAST ROW: YOU DO — alone, while the teacher circulates (17 min) ──────────
\mainidea[You do, alone]{Divide \& Factor} &
Four problems, each a step harder. Before every division, run the checklist.
\begin{center}
\begin{tikzpicture}[scale=1.0]
  \node[draw=plum, very thick, fill=lilac, rounded corners=3pt, inner sep=6pt,
        align=left, font=\small] at (0,0) {%
    \stepnum{1}\ Every power, highest to lowest --- a $0$ for each missing one.\\[3pt]
    \stepnum{2}\ Divisor $x - c$: use $c$. \ Divisor $x + c$: use $-c$.\\[3pt]
    \stepnum{3}\ The last number is the remainder, and it equals $p(c)$.};
\end{tikzpicture}
\end{center}
\\
& \notesprompt{On your own. Show your work.}
\begin{probgrid}{|Y|Y|}
\pcell{15}{$(2x^3 + 3x^2 - 4x + 5) \div (x + 2)$ by synthetic division. Write dividend $=$ divisor $\cdot$ quotient $+$ remainder.\par\smallskip
{\setlength{\tabcolsep}{2pt}\begin{tabular}{@{}>{\dstrut}K|KKKK@{}}
  &  &  &  &  \\
     &     &   &   &  \\ \cline{2-5}
     &  &  &  &  \\
\end{tabular}}}{1.3cm}{} &
\pcell{16}{$p(x) = x^4 - 3x^2 + 5$. Find the remainder when $p(x)$ is divided by $x - 2$ by synthetic division, then by finding $p(2)$.\par\smallskip
{\setlength{\tabcolsep}{2pt}\begin{tabular}{@{}>{\dstrut}K|KKKKK@{}}
  &  &  &  &  &  \\
     &     &   &   &  &  \\ \cline{2-6}
     &  &  &  &  &  \\
\end{tabular}}}{1.3cm}{} \\ \hline
\end{probgrid}
\\
& \begin{probgrid}*{|Y|Y|}
\pcell{17}{Show that $x - 4$ is a factor of $x^3 - 13x - 12$. Then find all the zeros.\par\smallskip
{\setlength{\tabcolsep}{2pt}\begin{tabular}{@{}>{\dstrut}K|KKKK@{}}
  &  &  &  &  \\
     &     &   &   &  \\ \cline{2-5}
     &  &  &  &  \\
\end{tabular}}}{1.6cm}{} &
\pcell{18}{Find the value of $k$ that makes $x + 2$ a factor of $x^3 + kx^2 - x - 6$. Explain using the Factor Theorem.}{2.8cm}{} \\ \hline
\end{probgrid}
\\ \hline
\end{guidednotes}

% ── THE NOTES END AT THE YOU DO ROW; AP Practice follows; homework is DeltaMath ─

\end{document}
